\documentclass[11pt]{article}
\usepackage[a4paper,margin=27mm]{geometry}
\usepackage{amsmath,amssymb,amsthm,hyperref}
\newtheorem{theorem}{Theorem}
\newtheorem{lemma}[theorem]{Lemma}
\newtheorem{proposition}[theorem]{Proposition}
\newcommand{\wt}{\operatorname{wt}}
\title{An exponential-size positive splitting in the reduced algebra}
\author{Walsh symmetrization and the full-activity obstruction}
\date{30 September 2026}
\begin{document}\maketitle
\noindent\textbf{Status.} Complete proof, independently reviewed by two collaborating agents.
No formal verification is claimed. The construction
produces rational piecewise constant fair densities, not finite uniform
dice: some cells have all letters active. The last proposition identifies
a genuine obstruction to eliminating those cells within this particular
family. No claim of literature novelty is made.

\section*{A multigraded cancellation lemma}
Work in the associative algebra over $\mathbb Q$ generated by
$H,Y_1,\ldots,Y_n$, truncated above total degree $n$, and quotient further
by every word in which some $Y_i$ occurs twice. There is no restriction
on repeated $H$. All exponentials and logarithms are finite formal sums.
For $b\in\mathbb F_2^n$, let $\rho_b$ fix $H$ and send
$Y_i$ to $(-1)^{b_i}Y_i$.

Choose rows $b_1,\ldots,b_L\in\mathbb F_2^n$. Starting with
$W_0=\exp(H+\sum_iY_i)$, define
\[
 W_\ell=W_{\ell-1}\rho_{b_\ell}(W_{\ell-1}).
\]
For $S\subseteq[n]$ put
\[
 w_\ell(S)=\#\{j\le\ell:\ \sum_{i\in S}(b_j)_i=1\pmod2\}.
\]
The $Y$-support of a nonzero monomial is the set of its $Y_i$ letters;
it contains each index at most once.

\begin{lemma}\label{cancel}
Every term of $\log W_\ell$ with nonempty $Y$-support $S$ has total
degree at least $w_\ell(S)+1$. The terms with empty $Y$-support equal
$2^\ell H$.
\end{lemma}
\begin{proof}
At $\ell=0$ the assertion is immediate. Write $A=\log W_{\ell-1}$
and expand $\log(e^A e^{\rho_{b_\ell}A})$ by the finite BCH formula.
If $S$ has odd parity under $b_\ell$, its contribution from
$A+\rho_{b_\ell}A$ cancels. Otherwise its old degree bound is precisely
the required new bound.

Each remaining BCH term is a bracket of $k\ge2$ homogeneous components
with disjoint nonempty supports $S_j$, possibly together with components
of empty support. Empty-support components have degree one. By induction
the degree of the bracket is at least
\[
 k+\sum_j w_{\ell-1}(S_j)
 \ge k+w_{\ell-1}\Bigl(\bigcup_jS_j\Bigr),
\]
since Hamming weight of a sum modulo two is at most the sum of the
Hamming weights. This is at least the old bound plus one, which is
sufficient whether the new parity is odd or even. Brackets containing
only $H$ vanish, proving the empty-support assertion.
\end{proof}

\begin{theorem}\label{path}
For each $n\ge1$ there is a positive piecewise constant path with
$2^{O(n)}$ equal rational cells, each velocity of the form
\[
 2\sum_{i\in S}X_i\qquad(S\subseteq[n]),
\]
whose signature in the reduced free algebra is exactly
$\exp(2^L\sum_iX_i)$. Equivalently, after normalization it is a
permutation-fair family of $n$ continuous probability distributions with
densities taking only the values zero and two.
\end{theorem}
\begin{proof}
Choose an injective binary linear code $B:\mathbb F_2^n\to
\mathbb F_2^L$ with minimum nonzero weight at least $n$, where $L=16n$
suffices. Indeed, take a random binary $L\times n$ matrix. For each
nonzero input the output has the $\operatorname{Bin}(L,1/2)$ weight law;
Chernoff's bound gives
\[
 \Pr(\wt(Bs)<n)\le e^{-49n/16}.
\]
A union bound over fewer than $2^n$ nonzero inputs is less than one.
Thus such a matrix exists, and its distance also ensures injectivity.

Use its rows in Lemma~\ref{cancel}. Every nonempty support then has
$w_L(S)\ge n$, so all its log terms vanish after truncation. Therefore
$W_L=\exp(2^L H)$. Substitute
\[
 H=\sum_iX_i,\qquad Y_i=X_i
\]
in the reduced algebra on the original $X_i$. This substitution is
well-defined: a repeated $Y_i$ becomes a repeated original letter and
vanishes. Each of the $2^L$ exponential factors has velocity
$\sum_i(1+\sigma_i)X_i$, where $\sigma_i\in\{-1,1\}$.
Finally scale each cell duration by $2^{-L}$. The first-level coefficients
are then one, and the full reduced signature is $\exp(\sum_iX_i)$,
which is exactly uniformity of every distinct-letter order probability.
\end{proof}

\section*{Why this does not yet finiteize by smaller alphabets}
\begin{proposition}\label{full}
In the construction above, each of the $2^n$ activity subsets occurs
exactly $2^{L-n}$ times. The assertion remains true after arbitrary
fixed sign changes of the $Y_i$ and after reordering the $L$ binary
symmetrizations. In particular, a full-activity cell cannot be eliminated
by selecting an affine phase or by reordering those symmetrizations.
\end{proposition}
\begin{proof}
The exponential factors are indexed by $t\in\mathbb F_2^L$, and their
sign vector is $((-1)^{(B^Tt)_i})_{i=1}^n$, in some binary order.
The map $B^T$ is onto because $B$ has rank $n$. Every sign vector has
exactly $2^{L-n}$ preimages. Fixed phases translate the sign vectors;
row reordering merely permutes the domain coordinates.
\end{proof}

The full-activity factor is a positive multiple of the uniform direction
$H$. At an endpoint of an already fair path it can be cancelled while
preserving fairness. Interior factors cannot in general be removed:
$A\exp(aH)B=\exp(cH)$ does not imply $AB=\exp((c-a)H)$.
Thus the present proof does not give a finite uniform-dice construction
of exponential size. Finiteizing every cell by a fair set on its active
alphabet is circular at the full-activity cells.

\begin{proposition}\label{half}
Suppose a family of nonconstant Walsh-character densities
$f_i=1+\epsilon_i\chi_i$, with $\epsilon_i\in\{-1,1\}$, has at most
$n/2$ active coordinates on every equal cell. If $n\ge4$, that family
cannot be three-wise fair.
\end{proposition}
\begin{proof}
Each nonconstant character has mean zero, so the mean number of active
coordinates is $n/2$. The assumed upper bound is therefore equality
on every cell, and $\sum_i\epsilon_i\chi_i=0$ identically. Independence
of distinct Walsh characters implies that each character occurs equally
often with the two signs.

Write $G_\chi(t)=\int_0^t\chi(s)\,ds$. The distinct CDF profiles are
$t+G_\chi(t)$ and $t-G_\chi(t)$. If two different character classes
occur, their four distinct CDFs obey the nontrivial linear relation
$(t+G_\chi)+(t-G_\chi)=(t+G_\psi)+(t-G_\psi)$. This contradicts the
proved theorem that distinct CDFs in a continuous three-wise fair family
are linearly independent. If only one class occurs, each of its two
distinct profiles occurs $n/2\ge2$ times, contradicting the companion
statement that at most one distinct CDF profile can be repeated.
Both inputs are proved in
\texttt{two\_dimensional\_equalization\_barrier.tex}.
\end{proof}
\section*{Related primary literature}
Walsh modulation and concatenated projection sequences are established
in dynamical decoupling. Qi, Dowling, and Viola identify their
correspondence and express cancellation order through Hamming weights
in their qubit setting; see
\href{https://arxiv.org/abs/1702.05533}{arXiv:1702.05533}, especially
pp.~2--3, published in \emph{Quantum Information Processing} 16 (2017),
article 272. The earlier paper by Hayes, Khodjasteh, Viola, and Biercuk,
\href{https://arxiv.org/abs/1109.6002}{arXiv:1109.6002}, develops Walsh
modulation for digital error suppression. The present proof specializes
the projection idea to a quotient in which repeated error labels vanish,
then uses an error-correcting code and a positive-density substitution.
The cited papers do not by themselves establish the dice construction
or remove the full-activity cells.
\end{document}
